\DocumentMetadata{
  pdfversion=2.0,
  pdfstandard=ua-2,
  lang=en-US,
  tagging=on,
  tagging-setup={math/setup=mathml-SE,tabsorder=structure}
}
\documentclass[11pt,letterpaper]{article}
\usepackage[margin=0.68in,includefoot]{geometry}
\usepackage{newcomputermodern}
\usepackage{amsmath,amscd,mathtools}
\usepackage{tikz,adjustbox}
\providecommand{\square}{\mdlgwhtsquare}
\usepackage[hidelinks]{hyperref}
\hypersetup{
  pdftitle={Topology PhD Exam, August 2025},
  pdfauthor={University of Florida Department of Mathematics},
  pdfsubject={Graduate mathematics examination},
  pdfdisplaydoctitle=true,
  bookmarksopen=true
}
\setcounter{secnumdepth}{0}
\setlength{\parindent}{0pt}
\setlength{\parskip}{0.28em}
\setlength{\emergencystretch}{3em}
\setlength{\leftmargini}{2.15em}
\setlength{\leftmarginii}{2.65em}
\setlength{\leftmarginiii}{3em}
\pagestyle{empty}
\raggedbottom
\allowdisplaybreaks
\NewTaggingSocketPlug{block/list/label}{examproblem}{%
  \tagstructbegin{tag=Lbl}\tagstructend%
  \tagstructbegin{tag=\LBody}%
  \tagstructbegin{tag=\ProblemHeadingTag,title-o={\ProblemHeadingText}}%
  \tagmcbegin{tag=\ProblemHeadingTag,actualtext={\ProblemHeadingText}}%
  #2%
  \tagmcend%
  \tagstructend%
}
\newenvironment{examproblems}[2]{%
  \def\ProblemHeadingTag{#2}%
  \AssignTaggingSocketPlug{block/list/label}{examproblem}%
  \begin{enumerate}%
  \setcounter{enumi}{#1}%
  \renewcommand{\labelenumi}{\textbf{\theenumi.}}%
  \setlength{\topsep}{0.35em}%
  \setlength{\partopsep}{0pt}%
  \setlength{\itemsep}{0.48em}%
  \setlength{\parsep}{0.18em}%
}{\end{enumerate}}
\newenvironment{examsubparts}{%
  \AssignTaggingSocketPlug{block/list/label}{default}%
  \begin{enumerate}%
  \setlength{\topsep}{0.18em}%
  \setlength{\partopsep}{0pt}%
  \setlength{\itemsep}{0.16em}%
  \setlength{\parsep}{0.1em}%
}{\end{enumerate}}
\newenvironment{examinstructions}{%
  \par\begingroup\itshape
}{%
  \par\endgroup\vspace{0.18em}
}
\makeatletter
\renewcommand\subsection{\@startsection{subsection}{2}{0pt}%
  {-1.25ex plus -.25ex minus -.1ex}%
  {0.55ex plus .1ex}%
  {\normalfont\large\bfseries}}
\renewcommand{\@maketitle}{%
  \newpage\null\vskip -1em%
  {\footnotesize\bfseries
    \href{https://www.ufl.edu/}{University of Florida}, \href{https://math.ufl.edu/}{Department of Mathematics}\par}%
  \vskip -0.5em%
  \tagstructbegin{tag=H1,title={Topology PhD Exam, August 2025}}%
  \tagpdfparaOff%
  \tagmcbegin{tag=H1}%
    {\LARGE\bfseries\href{https://gma.math.ufl.edu/exams/topology-phd/}{Topology PhD Exam}, August 2025\par}%
  \tagmcend%
  \tagpdfparaOn%
  \tagstructend%
  \vskip 0.25em%
  \tagmcbegin{artifact}%
    \hrule height 1pt%
  \tagmcend%
  \vskip 1em%
}
\makeatother
\title{Topology PhD Exam, August 2025}
\author{}
\date{}
\begin{document}
\maketitle
\thispagestyle{empty}
\begin{examinstructions}
Be neat and do not write too small.
\end{examinstructions}
\begin{examinstructions}
For the first five problems, show all of your work and support all statements.
\end{examinstructions}
\begin{examproblems}{0}{H2}
\def\ProblemHeadingText{Problem 1.}
\item
Let \((X,d)\) be a compact metric space. Let \(f:X\to X\) and suppose there is some \(c<1\) with \(d(f(x),f(y))\le c\cdot d(x,y)\) for all \(x,y\in X\). Prove there exists a unique \(x\in X\) with \(f(x)=x\).
\def\ProblemHeadingText{Problem 2.}
\item
Pick a CW complex structure on the Klein bottle~\(K\), and then use cellular homology to compute \(H_i(K;\mathbb{Z})\) for \(i=0,1,2\).
\def\ProblemHeadingText{Problem 3.}
\item
This problem has two parts.
\begin{examsubparts}
\item[\textnormal{(a)}]
Show that the antipodal map \(-1:S^n\to S^n\) defined by \(-1(x)=-x\) has degree \((-1)^{n+1}\).
\item[\textnormal{(b)}]
Prove that \(S^n\) has a continuous nonzero tangent vector field if and only if~\(n\) is odd.
\end{examsubparts}
\def\ProblemHeadingText{Problem 4.}
\item
Let~\(M\) be a compact orientable~\(n\)-manifold without boundary. Prove that if~\(n\) is odd, then the Euler characteristic \(\chi(M)\) is zero.
\def\ProblemHeadingText{Problem 5.}
\item
Let \(M_1\) and \(M_2\) be closed connected~\(n\)-manifolds, and let \(M_1\#M_2\) be their connected sum obtained by removing a ball from each and gluing along the resulting boundary spheres. Show that there are isomorphisms \(H_i(M_1\#M_2;\mathbb{Z})\cong H_i(M_1;\mathbb{Z})\oplus H_i(M_2;\mathbb{Z})\) for \(0<i<n\), with one exception: If both \(M_1\) and \(M_2\) are nonorientable, then \(H_{n-1}(M_1\#M_2;\mathbb{Z})\) is obtained from \(H_{n-1}(M_1;\mathbb{Z})\oplus H_{n-1}(M_2;\mathbb{Z})\) by replacing one of the two \(\mathbb{Z}_2\) summands by a~\(\mathbb{Z}\) summand.
\end{examproblems}
\begin{examinstructions}
Answer the following ten problems with complete definitions, complete statements, an example, or a short proof.
\end{examinstructions}
\begin{examproblems}{5}{H2}
\def\ProblemHeadingText{Problem 6.}
\item
Show that if \(g:S^2\to S^2\) is continuous and \(g(x)\ne g(-x)\) for all \(x\in S^2\), then~\(g\) is surjective.
\def\ProblemHeadingText{Problem 7.}
\item
State the Tychonoff Theorem.
\def\ProblemHeadingText{Problem 8.}
\item
Show that \(\mathbb{R}^{\omega}\) with the box topology is not connected.
\def\ProblemHeadingText{Problem 9.}
\item
Let \(N_g\) be the nonorientable surface of genus~\(g\), i.e., the~\(g\)-fold connected sum of projective planes. Use the Seifert–van Kampen theorem to derive the fundamental group of \(N_g\) for \(g\ge 1\).
\def\ProblemHeadingText{Problem 10.}
\item
Let~\(X\) be two unlinked circles in \(\mathbb{R}^3\) and let~\(Y\) be two simply linked circles. Show that the complements \(\mathbb{R}^3\setminus X\) and \(\mathbb{R}^3\setminus Y\) are not homotopy equivalent.
\def\ProblemHeadingText{Problem 11.}
\item
Show that if~\(n\) is even, then every map \(f:\mathbb{RP}^n\to\mathbb{RP}^n\) has a fixed point.
\def\ProblemHeadingText{Problem 12.}
\item
Let \(p_1:(\widetilde{X}_1,\widetilde{x}_1)\to(X,x_0)\) and \(p_2:(\widetilde{X}_2,\widetilde{x}_2)\to(X,x_0)\) be covering space maps, where all of the involved spaces are path-connected and locally path-connected. An isomorphism of covering spaces \(f:(\widetilde{X}_1,\widetilde{x}_1)\to(\widetilde{X}_2,\widetilde{x}_2)\) is a homeomorphism satisfying \(p_1=p_2f\). Prove the covering spaces \(p_1\) and \(p_2\) are isomorphic if and only if \(p_{1*}(\pi_1(\widetilde{X}_1,\widetilde{x}_1))=p_{2*}(\pi_1(\widetilde{X}_2,\widetilde{x}_2))\).
\def\ProblemHeadingText{Problem 13.}
\item
Define what it means for a subspace \(A\subseteq X\) to be a retract. Give an example of a retract \(A\subseteq X\) that is not a deformation retract.
\def\ProblemHeadingText{Problem 14.}
\item
Use the Tietze extension theorem to prove Urysohn’s lemma.
\def\ProblemHeadingText{Problem 15.}
\item
What are the cohomology groups and cup product structure in \(H^*((S^2\times S^8)\#(S^4\times S^6);\mathbb{Z})\), where \(\#\) denotes the connected sum of the manifolds? You do not need to prove your answer.
\end{examproblems}
\end{document}
